\section{An elementary certificate for the filtering moments}\label{app:numerical}

This appendix proves Proposition~\ref{prop:filtering-moments}.  There
are three errors to control.  Rational interval arithmetic encloses the
values computed on a 64-point grid.  The inner quadrature error compares
those values with the true score functions at the grid points.  The
outer quadrature error then compares the average of the true node
integrands with their circle integrals.  We treat these errors in that
order, after reducing the scores to two convolutions.
All integrals are normalized circle integrals, and all decimal
constants denote the indicated exact rationals.

\subsection{Reducing the scores to two convolutions}

Use the continuous functions of Section~\ref{sec:filtering}.  Define
\begin{equation}\label{eq:numerical-convolutions}
\begin{split}
 \mathcal P(x)&=\int F(x+w)G(w)\,dw,\\
 \mathcal Y(x)&=\int F(x+w)G(w)(r(x+w)+s(w))\,dw,\\
 L(x)&=\log Q(x),\qquad
 B(x)=\frac{\mathcal Y(x)}{\mathcal P(x)}-\log(Z_fZ_g),\\
 Z_{f,\kappa}&=\int F^\kappa,\qquad Z_{g,\kappa}=\int G^\kappa,\\
 H&=\log Z_{f,\kappa}+\log Z_{g,\kappa}
          -\kappa\log(Z_fZ_g).
\end{split}
\end{equation}
Here $Q=\mathcal P/(Z_fZ_g)$, and direct substitution in the score
definitions gives
\begin{equation}\label{eq:numerical-score-identities}
 D=B-L,\qquad
 u=L+(\kappa-1)B-H,\qquad
 v=(2-\kappa)B-2L+H.
\end{equation}
Indeed, $B$ is the $K_x$-mean of $\log(f(x+\cdot)g)$, whereas
$A+E=\kappa B-H$.  These identities reduce the computation to the two
convolutions in \eqref{eq:numerical-convolutions}.

Write $x_j=2\pi j/64$ and
$T_{64}(h)=64^{-1}\sum_{j=0}^{63}h(x_j)$.  A hat means that each
integral defining a quantity has been replaced by $T_{64}$; nonlinear
operations are performed after these replacements.  In particular
\[
 \widehat{\mathcal P}_j=\frac1{64}\sum_{i=0}^{63}F(x_{i+j})G(x_i),
 \quad
 \widehat{\mathcal Y}_j=\frac1{64}\sum_{i=0}^{63}
 F(x_{i+j})G(x_i)(r(x_{i+j})+s(x_i)),
\]
where indices are taken modulo 64.  Let $\widehat X_j$ and
$\widehat X'_j$ be the values of the first and second integrands in
\eqref{eq:continuous-filtering-moments}, respectively, obtained from
$\widehat L,\widehat B,\widehat H$ using
\eqref{eq:numerical-score-identities}.

\subsection{Finite arithmetic certificate}

Put $\beta=0.19$, and let $r_\beta,s_\beta$ be obtained from $r,s$ by
multiplying their $j$th cosine coefficients by $\cosh(\beta j)$.
The following enclosures will suffice.  In the first four rows the
error is at most $5\cdot10^{-7}$, and in the last two it is at most
$5\cdot10^{-6}$.
\begin{center}
\begin{tabular}{lc}
\toprule
Quantity & Center of enclosure\\
\midrule
$\widehat Z_f$ & $1.1963625$\\
$\widehat Z_g$ & $1.2877407$\\
$\widehat Z_{f,\kappa}$ & $1.2110790$\\
$\widehat Z_{g,\kappa}$ & $1.3038111$\\
$T_{64}(2^{2r_\beta})$ & $2.4901514$\\
$T_{64}(2^{2s_\beta})$ & $2.2607200$\\
\bottomrule
\end{tabular}
\end{center}

Table~\ref{tab:filter-certificate} gives the remaining node enclosures.
The errors in its four numerical columns are at most $4\cdot10^{-6}$,
$2\cdot10^{-5}$, $0.0004$, and $0.0004$, respectively.  All columns
are invariant under $j\mapsto64-j$.  Consequently, averaging the last
two columns with weight one at $j=0,32$ and weight two otherwise,
including their stated errors, gives
\begin{equation}\label{eq:hatted-average-bound}
 \frac1{64}\sum_{j=0}^{63}\widehat X_j<1.4295,
 \qquad
 \frac1{64}\sum_{j=0}^{63}\widehat X'_j<1.4295.
\end{equation}

\begin{table}[p]
\centering
\small
\begin{tabular}{r@{\qquad}rrrr}
\toprule
$j$ & $\widehat{\mathcal P}_j$ & $\widehat{\mathcal Y}_j$
& $\widehat X_j$ & $\widehat X'_j$\\
\midrule
0 &1.7255739&2.3814527&1.71479&1.72487\\
1 &1.7267975&2.3812891&1.70072&1.70914\\
2 &1.7298232&2.3798409&1.66451&1.66858\\
3 &1.7329951&2.3749116&1.62068&1.61932\\
4 &1.7343372&2.3645517&1.58491&1.57884\\
5 &1.7323079&2.3482031&1.56813&1.55943\\
6 &1.7263040&2.3267607&1.57405&1.56556\\
7 &1.7167948&2.3018438&1.59915&1.59372\\
8 &1.7051334&2.2751159&1.63351&1.63323\\
9 &1.6931319&2.2480097&1.66273&1.66812\\
10 &1.6824618&2.2214392&1.67219&1.68188\\
11 &1.6740048&2.1950779&1.65371&1.66486\\
12 &1.6673951&2.1665832&1.61070&1.62013\\
13 &1.6610380&2.1317550&1.55729&1.56268\\
14 &1.6527012&2.0862092&1.51133&1.51181\\
15 &1.6404589&2.0278888&1.48634&1.48247\\
16 &1.6235129&1.9586974&1.48675&1.48020\\
17 &1.6024686&1.8838599&1.50690&1.50018\\
18 &1.5789685&1.8092831&1.53309&1.52883\\
19 &1.5549462&1.7387429&1.54784&1.54794\\
20 &1.5318669&1.6725904&1.53630&1.54099\\
21 &1.5101580&1.6080342&1.49268&1.50055\\
22 &1.4888661&1.5397614&1.42405&1.43289\\
23 &1.4655935&1.4601151&1.34881&1.35678\\
24 &1.4368999&1.3596108&1.29024&1.29629\\
25 &1.3993369&1.2292982&1.26827&1.27187\\
26 &1.3509552&1.0651814&1.29131&1.29177\\
27 &1.2926991&0.8725981&1.34397&1.33998\\
28 &1.2289553&0.6674810&1.36465&1.35461\\
29 &1.1668804&0.4731120&1.23032&1.21474\\
30 &1.1147571&0.3140122&0.82842&0.81341\\
31 &1.0800632&0.2102374&0.27604&0.26987\\
32 &1.0678994&0.1742636&0&0\\
\bottomrule
\end{tabular}
\caption{Node enclosures for the filtering certificate.  The allowed
absolute errors are $4\cdot10^{-6}$, $2\cdot10^{-5}$, $0.0004$, and
$0.0004$, respectively.  Symmetry supplies the nodes $33,\ldots,63$.}
\label{tab:filter-certificate}
\end{table}

We specify elementary rational operations that reproduce these
enclosures.  The numbers $\cos x_j$, for $0\le j\le16$, lie within
$3\cdot10^{-10}$ of the following values, in order:
\[
\begin{array}{rrrrr}
1&0.9951847267&0.9807852804&0.9569403357&0.9238795325\\
0.8819212643&0.8314696123&0.7730104534&0.7071067812&0.6343932842\\
0.5555702330&0.4713967368&0.3826834324&0.2902846773&0.1950903220\\
0.0980171403&0
\end{array}
\]
For example, these intervals follow by using the cosine series through
degree 24, with remainder at most $|t|^{26}/26!$, and
\[
 3.1415926535897<\pi<3.1415926535899.
\]
The latter enclosure follows from
$\pi/4=\arctan(1/2)+\arctan(1/3)$ and the alternating arctangent
series through degree 159.  Reflection and a half-turn give the
remaining cosine intervals.  Taking the dot products with the
coefficients in \eqref{eq:filter-polynomials}, with indices reduced
modulo 64, gives rational intervals for every $r(x_j),s(x_j)$.

For all the exponential evaluations one may use
\begin{equation}\label{eq:certificate-exp}
 \exp(t)\in
 \sum_{i=0}^{45}\frac{t^i}{i!}
 +\left[-R_{\exp}(t),R_{\exp}(t)\right],\qquad
 R_{\exp}(t)=\frac{|t|^{46}}{46!\,(1-|t|/47)}.
\end{equation}
All arguments needed for the finite tables have $|t|\le8$.
For logarithms put $d=(y-1)/(y+1)$ and use
\begin{equation}\label{eq:certificate-log}
 \ln y\in
 2\sum_{i=0}^{15}\frac{d^{2i+1}}{2i+1}
 +\left[-R_{\log}(y),R_{\log}(y)\right],\qquad
 R_{\log}(y)=\frac{2|d|^{33}}{33(1-d^2)}.
\end{equation}
These follow by bounding the remaining absolutely convergent tails;
in particular they give
$0.693147180559<\ln2<0.693147180561$.
Compute powers of two as $\exp(t\ln2)$, and compute
$\cosh t=(\exp t+\exp(-t))/2$.
Apply interval addition, multiplication, and division to the displayed
finite sums and then to
\[
\begin{split}
 \widehat L_j&=\log\widehat{\mathcal P}_j
                    -\log\widehat Z_f-\log\widehat Z_g,\\
 \widehat B_j&=\widehat{\mathcal Y}_j/\widehat{\mathcal P}_j
                    -\log\widehat Z_f-\log\widehat Z_g,\\
 \widehat H&=\log\widehat Z_{f,\kappa}+\log\widehat Z_{g,\kappa}
                    -\kappa(\log\widehat Z_f+\log\widehat Z_g).
\end{split}
\]
These operations give the stated enclosures.  The intermediate
intervals may be kept at their rational endpoints; the displayed
decimal columns are final enclosures, not instructions to round
intermediate calculations.  Thus the tables and the two explicit
series remainder bounds constitute a finite arithmetic certificate.

\subsection{Controlling the inner quadratures}

We use the following form of the analytic trapezoid estimate
\cite[Theorem~3.2]{TrefethenWeideman2014}.  If a
$2\pi$-periodic function $h$ is analytic on a neighborhood of the
closed strip $|\operatorname{Im}z|\le\sigma$ and $|h(z)|\le B_0$
there, then
\begin{equation}\label{eq:analytic-trapezoid}
 \left|T_{64}(h)-\int h\right|
       \le\frac{2B_0}{e^{64\sigma}-1}.
\end{equation}
Indeed, shifting the contour for its $k$th Fourier coefficient toward
the appropriate boundary gives
$|\widehat h(k)|\le B_0e^{-|k|\sigma}$.  Averaging at the 64 nodes
retains precisely the coefficients indexed by multiples of 64; the
nonzero multiples have total absolute value at most the right side
of \eqref{eq:analytic-trapezoid}.

For the inner integrals defining the hatted values, take
$\sigma=0.45$.  The inequality
$|\cos(jz)|\le\cosh(j\sigma)$ and the displayed coefficients give
\[
 |r(z)|<6.50,\qquad |s(z)|<8.71,
 \qquad |2r_\beta(z)|<18.4,\qquad |2s_\beta(z)|<22.4.
\]
For instance these are obtained by summing the absolute coefficients
times $\cosh(0.45j)$; the two last bounds also include the factor
$2\cosh(0.19j)$.  Formula~\eqref{eq:analytic-trapezoid} gives the
following errors, distinct from the finite-table rounding errors:
\begin{equation}\label{eq:inner-quadrature-errors}
\begin{array}{c|c}
\text{Quantity}&\text{Error caused by quadrature}\\\hline
Z_f,Z_g,Z_{f,\kappa},Z_{g,\kappa}&4\cdot10^{-10}\\
\mathcal P(x_j)&3\cdot10^{-8}\\
\mathcal Y(x_j)&5\cdot10^{-7}\\
\int 2^{2r_\beta},\ \int 2^{2s_\beta}&5\cdot10^{-6}
\end{array}
\end{equation}
In detail the respective strip bounds can be taken as
$2^{\kappa\,8.71}$, $2^{6.50+8.71}$,
$(6.50+8.71)2^{6.50+8.71}$, and $2^{22.4}$.
Substituting them in \eqref{eq:analytic-trapezoid} proves every row of
\eqref{eq:inner-quadrature-errors} by elementary rational bounds on
the exponential.

Using the ranges in Table~\ref{tab:filter-certificate} to propagate
these errors through \eqref{eq:numerical-convolutions} gives, at every
node,
\begin{equation}\label{eq:node-score-errors}
 |L-\widehat L|<10^{-7},\qquad
 |B-\widehat B|<2\cdot10^{-6},\qquad
 |H-\widehat H|<5\cdot10^{-9}.
\end{equation}
For example, division is safe because
$\widehat{\mathcal P}_j>1.06789$, and
$|\log x-\log y|\le |x-y|/(\min(x,y)\ln2)$ controls the logarithms.
Applying \eqref{eq:numerical-score-identities} to
\eqref{eq:node-score-errors} shows that either true node integrand
is at most $1.00001$ times its hatted value.  This also holds at
$x_{32}=\pi$, where both values are zero.  We have therefore proved
\begin{equation}\label{eq:outer-node-bound}
 T_{64}(M^2 2^{8.5(\Delta-u)})<1.4295(1.00001),\qquad
 T_{64}(M^2 2^{8.3(av-u)})<1.4295(1.00001).
\end{equation}
It remains to control the outer quadrature, whose integrands involve
the logarithm and quotient in \eqref{eq:numerical-convolutions}.

\subsection{An analytic strip for the outer integrands}

The first table and \eqref{eq:inner-quadrature-errors} give
\begin{equation}\label{eq:shifted-l2}
 \int|f(w+i\beta)|^2\,dw
     =Z_f^{-2}\int2^{2r_\beta(w)}\,dw<1.761,
 \qquad
 \int|g(w+i\beta)|^2\,dw<1.38.
\end{equation}
The same bounds hold for imaginary shifts of absolute value at most
$\beta$.  To see this directly, the functions are real on the real
circle, so Parseval's identity pairs their positive and negative
Fourier coefficients into nonnegative multiples of
$\cosh(2k y)$; their squared norms are thus nondecreasing in $|y|$.
Their shifted means remain one by contour shifting.  Hence the
corresponding squared norms for $f-1$ and $g-1$ are less than $0.761$
and $0.38$.

For $z=x+iy$ with $|y|\le0.38=2\beta$, periodic contour shifting
in \eqref{eq:filter-kernel} yields
\[
 Q(z)=\int f(x+w+iy/2)g(w-iy/2)\,dw.
\]
Subtracting the product of the shifted means and applying
Cauchy--Schwarz gives
\begin{equation}\label{eq:Q-strip}
 |Q(z)-1|\le\sqrt{0.761\cdot0.38}<0.54.
\end{equation}
In particular $Q$ has no zeros on this strip, and
$L=\log Q=(\ln Q)/\ln2$ is analytic there with the branch real on
the real axis.  Moreover
\[
 -\operatorname{Re}L\le-\log0.46<1.121.
\]
The numerator $QB$ is also analytic.  The same split contour gives,
for every $|y|\le2\beta$,
\begin{align*}
 Q(z)B(z)
  =\int &f(x+w+iy/2)g(w-iy/2)\\
       &\mathord{}\cdot
       \bigl(r(x+w+iy/2)+s(w-iy/2)-\log(Z_fZ_g)\bigr)\,dw.
\end{align*}
This uses the analytic expression $r+s-\log(Z_fZ_g)$ for the
logarithm of the product, not a new branch choice.  Each factor is
evaluated within the strip of half-width $\beta$, where the
coefficient sums give $|r|<2.96$ and $|s|<4.67$.  Consequently, on
the full strip of half-width $2\beta$,
\begin{equation}\label{eq:B-strip}
 |Q B|\le\sqrt{1.761\cdot1.38}
       \bigl(2.96+4.67+\log Z_f+\log Z_g\bigr)<13,
 \qquad |B|<13/0.46.
\end{equation}
The normalizer enclosures also give $H<0.016$.

Now \eqref{eq:numerical-score-identities},
\eqref{eq:Q-strip}, and \eqref{eq:B-strip} imply
\[
 8.5\operatorname{Re}(\Delta-u)
 \le8.5\left(0.02776+1.121+0.0315\frac{13}{0.46}+0.016\right)
 <17.6.
\]
For the second exponent use the exact identity
\[
 av-u=-(1+2a)L+
       \bigl(a(2-\kappa)-(\kappa-1)\bigr)B+(1+a)H.
\]
The same bounds give
\[
 8.3\operatorname{Re}(av-u)
 \le8.3\left((1+2a)1.121+
       |a(2-\kappa)-(\kappa-1)|\frac{13}{0.46}
       +(1+a)0.016\right)<17.6.
\]
Finally $|2+2\cos z|\le2+2\cosh(0.38)<4.16$.
Both outer integrands therefore extend analytically to the strip
$|\operatorname{Im}z|\le0.38$ and have absolute value below $2^{20}$
there.  (The strict bounds also give a neighborhood of the closed
strip.)  Formula~\eqref{eq:analytic-trapezoid} bounds either outer
quadrature error by
\[
 \frac{2^{21}}{e^{64\cdot0.38}-1}<0.00007.
\]
Combining this with \eqref{eq:outer-node-bound}, each integral in
\eqref{eq:continuous-filtering-moments} is strictly less than
\[
 1.4295(1.00001)+0.00007
   =1.429584295<2^{0.5168}.
\]
The last comparison follows already from
$2^{0.5168}>1.4307781$, obtained from
\eqref{eq:certificate-exp}.  This proves both filtering moments.
